\chapter{Reproducible Illustrative Computation}
\label{app:computation}

The project includes \texttt{scripts/illustrative\_model.py}. It generates the normalized figures and \texttt{illustrative\_results.csv}. The script requires Python, NumPy, SciPy, and Matplotlib.

\section{Execution}

From the project root:
\begin{lstlisting}[language=bash]
python scripts/illustrative_model.py
latexmk -pdf -interaction=nonstopmode main.tex
\end{lstlisting}

\section{Algorithm}

\begin{enumerate}
\item Construct a Maxwell--J\"uttner energy density on a fine dimensionless energy grid.
\item Normalize by numerical quadrature.
\item Select cutoff $E_c$ and set the distribution above it to zero.
\item Integrate the removed number fraction.
\item Add the removed number to a normalized Gaussian kernel below the cutoff.
\item Renormalize to remove quadrature error.
\item Calculate moments $M_q$ and ratios to the reference distribution.
\item Construct a surrogate continuum photon spectrum and smooth energy-dependent suppression.
\item Apply three response weightings to illustrate why source-power and dose-like ratios differ.
\end{enumerate}

\section{Replacement by a research-grade calculation}

A final calculation should replace the surrogate steps with:
\begin{enumerate}
\item a relativistic Fokker--Planck solution for $f(p,\xi)$;
\item exact or benchmarked electron--ion differential cross sections;
\item an electron--electron kernel integrated over two-particle phase space;
\item synchrotron spectral-angular emission;
\item collector electron transport and thick-target bremsstrahlung;
\item Monte Carlo transport through the actual CAD geometry; and
\item measured detector response and uncertainty covariance.
\end{enumerate}

\section{Suggested file interfaces}

\begin{longtable}{p{0.25\textwidth}p{0.65\textwidth}}
\toprule
File & Contents \\
\midrule
\texttt{distribution.h5} & grids in position, momentum, pitch angle, time; $f_e$ and uncertainty \\
\texttt{source.h5} & photon energy, angle, position, time, channel, and emitted weight \\
\texttt{transport.inp} & geometry, materials, physics settings, source reference \\
\texttt{dose.csv} & voxel or detector identifier, dose, statistical uncertainty \\
\texttt{biology.csv} & sample identifier, delivered dose, endpoint, time, replicate \\
\texttt{provenance.yaml} & software versions, git commit, calibration identifiers \\
\bottomrule
\end{longtable}
